\chapter{Conclusion and Future Work}
\label{ch:conclusion}

\section{Conclusion}
This thesis has proposed and validated \textbf{SENTINEL} - a two-tier cognitive architecture designed to solve the alert fatigue problem in Security Operations Centers (SOCs). By integrating a high-speed pre-filter (Tier-1) and a Machine Learning Gateway (LightGBM), the system successfully automated the resolution of $83.8\%$ of network logs in merely $0.35$\,ms. This innovation slashed the end-to-end processing latency by $82.97\%$ compared to standalone LLM deployments.
For complex, ambiguous threats, Tier-2 leverages a locally hosted Agentic AI (Gemma-9B) combined with hybrid Retrieval-Augmented Generation (Dual-RAG) and Threat Memory to execute highly precise mitigations without violating Zero Trust security principles. Furthermore, the implementation of robust cryptographic Guardrails completely shielded the LLM from log-based Prompt Injection attacks, proving the enterprise-grade safety of the solution.

\section{Future Work}
While the system has achieved promising results, several pathways remain to elevate SENTINEL to production-grade scalability.
First, the architecture must transition to a Microservices model to enable horizontal scale-out on distributed systems like Kubernetes, replacing Redis Streams with Kafka to handle massive bandwidths without memory exhaustion.
Second, the Tier-1 pre-filter could be upgraded with lightweight deep-learning models or re-implemented in system-level languages (Go/Rust) to push latency down to the microsecond threshold.
Finally, a Multi-agent network could be developed to expand the system's capabilities beyond mere threat detection, fully automating digital forensics and deep incident response playbooks.

\vspace{0.6cm}
\begin{flushright}
    \textit{Closing a humble endeavor, yet opening profound aspirations in the journey to secure the digital realm through the power of Agentic AI.}
\end{flushright}
